\columnbreak
\section{Who Must Attack}

\begin{rulebox}{General Rule}
Each non-disrupted unit exerts a Zone of Control (ZOC) into the six hexes surrounding the one it occupies. All of the active player's ordered units in the ZOC of an enemy unit at the beginning of the active player's combat phase must attack one or more enemy units during the combat phase. The active player must attack all enemy ordered units in a friendly ZOC. The active player's disrupted units cannot attack, but the inactive player's disrupted units can be attacked at the active player's option. Artillery pieces cannot attack or be attacked in the combat phase. Leaders that are not stacked with a friendly unit cannot attack or be attacked. Leaders stacked with a friendly unit add their command rating to the combat strength of the unit with which they are stacked.
\end{rulebox}

\ruleslabel{Procedure}

At the beginning of his combat phase, before resolving any combats, the active player announces which of his units are attacking and which enemy units each is attacking. Only when all combats have been announced is the first combat resolved.

\caseheading{9.1} All ordered units exert a Zone of Control (ZOC) on the six adjacent hexes at all times.

Leaders, artillery pieces, and disrupted units never exert a ZOC. ZOCs are unaffected by the presence of any type of terrain or playing piece. If an ordered unit is in the ZOC of an enemy unit, the enemy unit is also in the ZOC of that ordered unit. Friendly and enemy ZOCs coexist in the same hex. There is no additional effect from having more than one ZOC exerted on a hex.

\caseheading{9.2} The manner in which attacks are divided is entirely up to the active player. The active player chooses which of his units will attack which enemy units.

All of the inactive player's ordered units that are in the ZOC of one of the active player's units must be attacked, and all the active player's ordered units that are in an enemy unit's ZOC must make an attack. Within this restriction, the active player can divide his attacks in whatever way he feels will give him the best results.

\caseheading{9.3} A unit can only attack adjacent enemy units. A given enemy unit can be attacked in a single combat by up to six units, each occupying one of the six adjacent hexes.

\caseheading{9.4} A unit cannot attack or be attacked more than once per combat phase.

When more than one unit is attacking the same enemy unit, all attacking units must execute their attacks in a single combined combat. The combat strengths of all the attacking units are combined into a single total combat strength for purposes of this attack.

\caseheading{9.5} Each unit's combat strength is a single indivisible number. A unit's combat strength cannot be divided among different combats when either attacking or defending.

\caseheading{9.6} When the active player's attacking units in more than one hex are all adjacent to the same defending units in more than one hex, the attacking units can resolve their attacks as a single combined attack, at the active player's option.

All of the attacking units must be adjacent to all the defending units for this option to be exercised. In such cases, the combat strengths of the attacking units are combined into a single total combat strength, and the combat strengths of the defending units are combined into a single total combat strength.

\caseheading{9.7} Modifications to the combat strengths of individual units are made before the combat strengths are combined.

If the defending units occupy more than one type of terrain, the combat is resolved on the single terrain line on the Combat Table most favorable to the defending units.

\caseheading{9.8} The attacker can divide his attacks so that some attacks are executed at an unfavorable combat ratio in order to obtain a more favorable combat ratio in other attacks.

Combat ratios lower than the lowest combat ratio shown on the terrain line on which the combat is being resolved are treated as the lowest combat ratio shown on that line. Combat ratios higher than the highest combat ratio shown on the line are resolved at the highest combat ratio shown.
